A diferencia de las variables de origen, los campos de las Tablas~\ref{tab:stints_wnba} y \ref{tab:coincidencias_wnba} forman parte del diseño de los datos procesados del proyecto. Se pretende generarlos mediante cruces y operaciones sobre registros históricos; por tanto, su presencia no está garantizada en los archivos Parquet originales.

{\small
\renewcommand{\arraystretch}{1.15}
\begin{longtable}{@{}p{3.8cm}p{1.7cm}p{6.9cm}@{}}
\caption{Variables propuestas para identificar intervalos entre sustituciones.}\label{tab:stints_wnba}\\
\toprule \textbf{Variable} & \textbf{Tipo} & \textbf{Procedimiento de obtención}\\\midrule\endfirsthead
\multicolumn{3}{l}{\small Continuación de la Tabla~\ref{tab:stints_wnba}}\\
\toprule\textbf{Variable}&\textbf{Tipo}&\textbf{Procedimiento de obtención}\\\midrule\endhead
\midrule\multicolumn{3}{r}{\small Continúa en la siguiente página}\\\endfoot
\bottomrule\endlastfoot
\texttt{stint\_id}&Texto&Identificador generado para el intervalo en el que no cambia ninguno de los quintetos.\\
\texttt{start\_clock}&Tiempo&Tiempo restante cuando inicia el intervalo.\\
\texttt{end\_clock}&Tiempo&Tiempo restante al concluir el intervalo.\\
\texttt{stint\_seconds}&Decimal&Duración, en segundos, después de normalizar el reloj y el periodo.\\
\texttt{home\_team\_id}&ID&Equipo local del encuentro.\\
\texttt{away\_team\_id}&ID&Equipo visitante del encuentro.\\
\texttt{score\_margin\_start}&Entero&Diferencia entre marcadores al inicio.\\
\texttt{score\_margin\_end}&Entero&Diferencia entre marcadores al cierre.\\
\texttt{focal\_lineup}&Conjunto&Cinco identificadores del quinteto propio, cuando puedan reconstruirse.\\
\texttt{opponent\_lineup}&Conjunto&Cinco identificadores de las rivales, cuando puedan reconstruirse.\\
\end{longtable}
}

\noindent\textbf{Descripción de la tabla.}
La Tabla~\ref{tab:stints_wnba} define la estructura propuesta para analizar segmentos del partido durante los cuales se mantienen las mismas diez jugadoras. Esta unidad permitiría asociar marcador, duración y quintetos; sin embargo, su construcción depende de reconstruir correctamente sustituciones y otras interrupciones de juego. Los campos temporales y los identificadores de quinteto son derivados, no columnas documentadas como tales en \texttt{pbp} \cite{wehoopPbp}.

\medskip

{\small
\renewcommand{\arraystretch}{1.15}
\begin{longtable}{@{}p{3.8cm}p{1.7cm}p{6.9cm}@{}}
\caption{Variables propuestas para coincidencias contextuales entre jugadoras rivales.}\label{tab:coincidencias_wnba}\\
\toprule\textbf{Variable}&\textbf{Tipo}&\textbf{Procedimiento de obtención}\\\midrule\endfirsthead
\multicolumn{3}{l}{\small Continuación de la Tabla~\ref{tab:coincidencias_wnba}}\\
\toprule\textbf{Variable}&\textbf{Tipo}&\textbf{Procedimiento de obtención}\\\midrule\endhead
\midrule\multicolumn{3}{r}{\small Continúa en la siguiente página}\\\endfoot
\bottomrule\endlastfoot
\texttt{focal\_player\_id}&ID&Identificador de jugadora del equipo analizado presente en el intervalo.\\
\texttt{opponent\_player\_id}&ID&Identificador de jugadora rival presente simultáneamente.\\
\texttt{focal\_role}&Categoría&Perfil funcional Guard, Wing o Big asignado mediante el análisis.\\
\texttt{opponent\_role}&Categoría&Perfil de la jugadora rival.\\
\texttt{compatible\_roles}&Binario&Indicador de compatibilidad de perfiles según la regla que se valide.\\
\texttt{focal\_starter}&Binario&Indica si la jugadora comenzó el encuentro, cuando pueda comprobarse.\\
\texttt{opponent\_starter}&Binario&Titularidad de la jugadora rival.\\
\texttt{shared\_seconds}&Decimal&Segundos de coincidencia dentro del intervalo.\\
\texttt{shared\_possessions}&Entero&Posesiones reconstruidas mientras ambas están en cancha.\\
\texttt{focal\_off\_possessions}&Entero&Posesiones ofensivas propias durante esa coincidencia.\\
\texttt{focal\_def\_possessions}&Entero&Posesiones defensivas durante la coincidencia.\\
\texttt{shared\_seconds\_game}&Decimal&Suma de segundos de coincidencia en un encuentro.\\
\texttt{shared\_possessions\_game}&Entero&Suma de posesiones coincidentes en un encuentro.\\
\end{longtable}
}

\noindent\textbf{Descripción de la tabla.}
La Tabla~\ref{tab:coincidencias_wnba} especifica las observaciones necesarias para una comparación contextual individual: cuánto tiempo coincidieron dos rivales y qué ocurrió durante ese periodo. La presencia simultánea no implica que una jugadora defendiera directamente a la otra. Los perfiles y la compatibilidad también deben calcularse; las estadísticas por intervalo no podrán tomarse de los totales de temporada sin procesar.